\section{Introduction}
\label{sec:intro}

%Throughout the years, 
Multiview-based 3D scene reconstruction~\cite{schoenberger2016sfm, yao2018mvsnet, xie2019pix2vox, murez2020atlas, wang2021multi, sayed2022simplerecon, xu2023unifying, wang2024dust3r, duisterhof2024mast3r, pan2024glomap} and novel view synthesis~\cite{mildenhall2020nerf, yu2021pixelnerf, yu2022plenoxels, mueller2022instant, zheng2024gpsgaussian, gslrm2024, charatan2024pixelsplat, sabour2024spotlesssplats, flynn2024quark, chen2025mvsplat} have been central topics in computer vision. 
These methods leverage multiview information to reconstruct high-fidelity scenes. 
However, data acquisition remains a bottleneck. 
Most data is collected manually, which is time-consuming and labor-intensive, follows preplanned camera trajectories, or relies on non-active capture systems~\cite{xu2023vr, broxton2020immersive, lin2021deep, yoon2020novel, li2022neural, lu2024diva, chen2024x360}.
%that do not generalize well to diverse objects.
%shapes.

To reduce human effort, next-best-view (NBV) planning has been explored for active capture~\cite{monica2018contour, zha1997next, liu2018object, hardouin2020surface, hardouin2020next, lin2022active, pan2022activenerf, lee2022uncertainty, zhan2022activermap, jin2023neu, sunderhauf2023density, ran2023neurar, guedon2023macarons, jiang2023fisherrf}. 
Traditional next-best-view methods rely on heuristic rules that can work well in specific scenarios but often fail to transfer because fixed rules or hyperparameters do not adapt across scenes~\cite{chen2024gennbv}. 
Learning-based next-best-view planners, including online-learning and generalizable methods, improve over preplanned trajectories, which frequently miss occluded regions.
Within learning-based approaches, reinforcement learning-based (RL-based) generalizable methods~\cite{peralta2020next, chen2024gennbv}, which pretrain on a dataset to avoid online learning and directly predict viewpoints as actions, show promising results.
This removes candidate-viewpoint sampling, which may potentially miss the best views and slow viewpoint acquisition. 
An occupancy-grid formulation~\cite{chen2024gennbv} further demonstrates strong coverage, viewpoint flexibility, and generalization. 
Nevertheless, performance remains limited and insufficiently robust across diverse object geometries.

To address the shortcomings, we propose Hestia, a Voxel-Face-Aware \underline{H}ierarchical N\underline{e}xt-Be\underline{s}\underline{t}-V\underline{i}ew \underline{A}cquisition for Efficient 3D Reconstruction. 
Hestia actively collects data in object-centric scenes by predicting five-degree-of-freedom (5-DoF) viewpoints (x, y, z, yaw, pitch) from voxel-face observations.
Specifically, Hestia systematically defines the next-best-view task by proposing core components such as dataset choice, observation and reward design, action space, and learning schemes, forming a foundation for the planner.

\begin{figure}[t]  
  \centering
  \includegraphics[width=\linewidth]{figures/voxel_ray.png}
  \caption{\textbf{A voxel is worth more than a ray.} 
  Unlike the RL-based generalizable method~\cite{chen2024gennbv}, Hestia treats each voxel as a cube by considering its six faces, rather than a point. This reduces the information loss inherent in point approximations, ensuring a more accurate representation of the voxel.
  }
  \label{fig:voxel_ray}
\end{figure}

An idea is \underline{\textit{``A voxel is worth more than a ray''}}. 
We incorporate the visibility of the six faces of each voxel into both the observation and the reward function (see~\cref{fig:voxel_ray,sec:methods}). 
Theoretically, if we sample with a one-ray camera in a scene with \(k\) unit cubes and treat each voxel as a point, then from a coupon collector’s perspective~\cite{boneh1997coupon} approximately \(k^{-1/6}\) of the faces will be missed when sampling stops (see Sec.~\ref{asec:theory}). 
Treating each voxel as a cube enables full face coverage, so Hestia accounts for individual voxel-face visibility and captures data more comprehensively. 
This representation adds little computational overhead and still achieves real-time operation at 25 FPS (see~\cref{tab:avg_bench}).

Hestia further improves learning by refining the observation, action, and learning process. 
For observation, Hestia uses the largest dataset that we processed from Objaverse~\cite{Deitke_2023_CVPR, deitke2024objaverse} for the next-best-view task, exposing the policy to a broad range of surface geometries rather than mostly cubic shapes (see Sec.~S8). 
For action, instead of predicting the full 5-DoF next-best view in one step, Hestia adopts a hierarchical structure. 
The policy first predicts a look-at point as the target of attention, then determines the viewpoint position conditioned on this point. 
For learning, Hestia formulates the task as a close-greedy optimization problem in which, given an occupancy grid, it selects the view that maximizes the current coverage ratio. 
The policy relies only on the previous image, the previous camera pose, and the occupancy grid, rather than a long sequence of images and poses~\cite{peralta2020next, chen2024gennbv}. 
We also use a small reward discount factor \(\gamma\) to prioritize immediate improvements without depending on an oversized terminal reward. 
Similar to a greedy algorithm, this reduces spurious correlations\footnote{Spurious correlations~\cite{interpretation1971spurious, kim2024discovering} refer to certain groups contributing to model errors. In this study, spurious correlations refer to large positive future rewards assigned to suboptimal current next-best-view decisions, leading to ineffective policy learning.} between current actions and large future rewards (see Fig.~S8). 
As a result, Hestia reaches higher coverage with fewer images than prior methods~\cite{guedon2023macarons, jin2023neu, chen2024gennbv, peralta2020next}.
Finally, we demonstrate that Hestia, as a next-best-view planner, is feasible for real-world application using a drone with an RGB camera as a mobile agent and a depth predictor~\cite{duisterhof2024mast3r, wang2024dust3r} to convert RGB images into depth maps. 
The contributions of this work are as follows:
\begin{itemize}
\item A RL-based generalizable next-best-view planner that considers voxels as cubes rather than points to avoid geometry overlooking.
\item A hierarchical structure for handling the high-dimensional continuous action space, a larger and more diverse training set for promoting robustness, and a close-greedy strategy for reducing spurious correlations.
\item Comprehensive evaluations on three datasets show that Hestia achieves non-marginal improvements and is suitable for 3D reconstruction under limited acquisition budgets.
\end{itemize}

